% OpenStax Calculus Volume 3 -- Section 4.6 Exercises: Directional Derivatives and the Gradient
% Exercises reproduced from OpenStax, licensed under CC BY 4.0.
% Source: https://openstax.org/books/calculus-volume-3/pages/4-6-directional-derivatives-and-the-gradient
\documentclass{exercises}
\title{OpenStax Calculus Volume 3}
\subtitle{Section 4.6 Exercises: Directional Derivatives and the Gradient}
\sourceurl{https://openstax.org/books/calculus-volume-3/pages/4-6-directional-derivatives-and-the-gradient}
\author{}
\date{}

\begin{document}
\maketitle


For the following exercises, find the directional derivative using the limit definition only.

\begin{enumerate}[start=1]
\item $f(x,y)=5-2x^{2}-\frac{1}{2}y^{2}$ at point $P(3,4)$ in the direction of $\mathbf{u}=(\cos \frac{\pi}{4})\mathbf{i}+(\sin \frac{\pi}{4})\mathbf{j}$
\item $f(x,y)=y^{2}\cos(2x)$ at point $P(\frac{\pi}{3},2)$ in the direction of $\mathbf{u}=(\cos \frac{\pi}{4})\mathbf{i}+(\sin \frac{\pi}{4})\mathbf{j}$
\item Find the directional derivative of $f(x,y)=y^{2}\sin(2x)$ at point $P(\frac{\pi}{4},2)$ in the direction of $\mathbf{u}=5\mathbf{i}+12\mathbf{j}$.
\end{enumerate}

For the following exercises, find the directional derivative of the function at point $P$ in the direction of $\mathbf{u}$ or $\mathbf{v}$ as appropriate.

\begin{enumerate}[resume]
\item $f(x,y)=xy$, $P(0,-2)$, $\mathbf{v}=\frac{1}{2}\mathbf{i}+\frac{\sqrt{3}}{2}\mathbf{j}$
\item $h(x,y)=e^{x}\sin y,P(1,\frac{\pi}{2}),\mathbf{v}=-\mathbf{i}$
\item $h(x,y,z)=xyz,P(2,1,1),\mathbf{v}=2\mathbf{i}+\mathbf{j}-\mathbf{k}$
\item $f(x,y)=xy,P(1,1),\mathbf{u}=\langle \frac{\sqrt{2}}{2},\frac{\sqrt{2}}{2}\rangle$
\item $f(x,y)=x^{2}-y^{2},\begin{matrix}\mathbf{u}=\langle \frac{\sqrt{3}}{2},\frac{1}{2}\rangle, & P(1,0)\end{matrix}$
\item $f(x,y)=3x+4y+7,\begin{matrix}\mathbf{u}=\langle \frac{3}{5},\frac{4}{5}\rangle, & P(0,\frac{\pi}{2})\end{matrix}$
\item $\begin{matrix}f(x,y)=e^{x}\cos y, & \mathbf{u}=\langle0,1\rangle, & P=(0,\frac{\pi}{2})\end{matrix}$
\item $\begin{matrix}f(x,y)=y^{10}, & \mathbf{u}=\langle0,-1\rangle, & P=(1,-1)\end{matrix}$
\item $f(x,y)=\ln(x^{2}+y^{2}),\begin{matrix}\mathbf{u}=\langle \frac{3}{5},\frac{4}{5}\rangle, & P(1,2)\end{matrix}$
\item $f(x,y)=x^{2}y,\begin{matrix}P(-5,5), & \mathbf{v}=3\mathbf{i}-4\mathbf{j}\end{matrix}$
\item $f(x,y,z)=y^{2}+xz,\begin{matrix}P(1,2,2), & \mathbf{v}=\langle2,-1,2\rangle\end{matrix}$
\end{enumerate}

For the following exercises, find the directional derivative of the function in the direction of the unit vector $\mathbf{u}=\cos \theta \mathbf{i}+\sin \theta \mathbf{j}$.

\begin{enumerate}[resume]
\item $f(x,y)=x^{2}+2y^{2},\theta=\frac{\pi}{6}$
\item $f(x,y)=\frac{y}{x+2y},\theta=-\frac{\pi}{4}$
\item $f(x,y)=\cos(3x+y),\theta=\frac{\pi}{4}$
\item $w(x,y)=ye^{x},\theta=\frac{\pi}{3}$
\item $\begin{matrix}f(x,y)=x\arctan(y), & \theta=\frac{\pi}{2}\end{matrix}$
\item $\begin{matrix}f(x,y)=\ln(x+2y), & \theta=\frac{\pi}{3}\end{matrix}$
\end{enumerate}

For the following exercises, find the gradient.

\begin{enumerate}[resume]
\item Find the gradient of $f(x,y)=\frac{14-x^{2}-y^{2}}{3}$. Then, find the gradient at point $P(1,2)$.
\item Find the gradient of $f(x,y,z)=xy+yz+xz$ at point $P(1,2,3)$.
\item Find the gradient of $f(x,y,z)$ at $P$ and the directional derivative in the direction of $\mathbf{u}$: $f(x,y,z)=\ln(x^{2}+2y^{2}+3z^{2}),\begin{matrix}P(2,1,4), & \mathbf{u}=\frac{-3}{13}\mathbf{i}-\frac{4}{13}\mathbf{j}-\frac{12}{13}\mathbf{k}\end{matrix}$.
\item $f(x,y,z)=4x^{5}y^{2}z^{3},\begin{matrix}P(2,-1,1), & \mathbf{u}=\frac{1}{3}\mathbf{i}+\frac{2}{3}\mathbf{j}-\frac{2}{3}\mathbf{k}\end{matrix}$
\end{enumerate}

For the following exercises, find the directional derivative of the function at point $P$ in the direction of $Q$.

\begin{enumerate}[resume]
\item $f(x,y)=x^{2}+3y^{2},\begin{matrix}P(1,1), & Q(4,5)\end{matrix}$
\item $f(x,y,z)=\frac{y}{x+z},\begin{matrix}P(2,1,-1), & Q(-1,2,0)\end{matrix}$
\end{enumerate}

For the following exercises, find the derivative of the function at $P$ in the direction of $\mathbf{u}$.

\begin{enumerate}[resume]
\item $f(x,y)=-7x+2y,\begin{matrix}P(2,-4), & \mathbf{u}=4\mathbf{i}-3\mathbf{j}\end{matrix}$
\item $f(x,y)=\ln(5x+4y),\begin{matrix}P(3,9), & \mathbf{u}=6\mathbf{i}+8\mathbf{j}\end{matrix}$
\item \techrequired Use technology to sketch the level curve of $f(x,y)=4x-2y+3$ that passes through $P(1,2)$ and draw the gradient vector at $P$.
\item \techrequired Use technology to sketch the level curve of $f(x,y)=x^{2}+4y^{2}$ that passes through $P(-2,0)$ and draw the gradient vector at $P$.
\end{enumerate}

For the following exercises, find the gradient vector at the indicated point.

\begin{enumerate}[resume]
\item $f(x,y)=xy^{2}-yx^{2},P(-1,1)$
\item $f(x,y)=xe^{y}-\ln(x),P(3,0)$
\item $f(x,y,z)=xy-\ln(z),P(2,-2,2)$
\item $f(x,y,z)=x\sqrt{y^{2}+z^{2}},P(-2,-1,-1)$
\end{enumerate}

For the following exercises, find the derivative of the function.

\begin{enumerate}[resume]
\item $f(x,y)=x^{2}+xy+y^{2}$ at point $(-5,-4)$ in the direction the function increases most rapidly
\item $f(x,y)=e^{xy}$ at point $(6,7)$ in the direction the function increases most rapidly
\item $f(x,y)=\arctan(\frac{y}{x})$ at point $(-9,9)$ in the direction the function increases most rapidly
\item $f(x,y,z)=\ln(xy+yz+zx)$ at point $(-9,-18,-27)$ in the direction the function increases most rapidly
\item $f(x,y,z)=\frac{x}{y}+\frac{y}{z}+\frac{z}{x}$ at point $(5,-5,5)$ in the direction the function increases most rapidly
\end{enumerate}

For the following exercises, find the maximum rate of change of $f$ at the given point and the direction in which it occurs.

\begin{enumerate}[resume]
\item $f(x,y)=xe^{-y}$, $(1,0)$
\item $f(x,y)=\sqrt{x^{2}+2y}$, $(4,10)$
\item $f(x,y)=\cos(3x+2y),(\frac{\pi}{6},-\frac{\pi}{8})$
\end{enumerate}

For the following exercises, find equations of


(a) the tangent plane and


(b) the normal line to the given surface at the given point.

\begin{enumerate}[resume]
\item The level surface $f(x,y,z)=12$ for $f(x,y,z)=4x^{2}-2y^{2}+z^{2}$ at point $(2,2,2)$.
\item $f(x,y,z)=xy+yz+xz=3$ at point $(1,1,1)$
\item $f(x,y,z)=xyz=6$ at point $(1,2,3)$
\item $f(x,y,z)=xe^{y}\cos z-z=1$ at point $(1,0,0)$
\end{enumerate}

For the following exercises, solve the problem.

\begin{enumerate}[resume]
\item The temperature $T$ in a metal sphere is inversely proportional to the distance from the center of the sphere (the origin: $(0,0,0)$). The temperature at point $(1,2,2)$ is $120^\circ \text{C}$.
\begin{enumerate}[label=(\alph*)]
  \item Find the rate of change of the temperature at point $(1,2,2)$ in the direction toward point $(2,1,3)$.
  \item Show that, at any point in the sphere, the direction of greatest increase in temperature is given by a vector that points toward the origin.
\end{enumerate}
\item The electrical potential (voltage) in a certain region of space is given by the function $V(x,y,z)=5x^{2}-3xy+xyz$.
\begin{enumerate}[label=(\alph*)]
  \item Find the rate of change of the voltage at point $(3,4,5)$ in the direction of the vector $\langle1,1,-1\rangle$.
  \item In which direction does the voltage change most rapidly at point $(3,4,5)$?
  \item What is the maximum rate of change of the voltage at point $(3,4,5)$?
\end{enumerate}
\item If the electric potential at a point $(x,y)$ in the $xy$-plane is $V(x,y)=e^{-2x}\cos(2y)$, then the electric intensity vector at $(x,y)$ is $\mathbf{E}=-\nabla V(x,y)$.
\begin{enumerate}[label=(\alph*)]
  \item Find the electric intensity vector at $(\frac{\pi}{4},0)$.
  \item Show that, at each point in the plane, the electric potential decreases most rapidly in the direction of the vector $\mathbf{E}$.
\end{enumerate}
\item In two dimensions, the motion of an ideal fluid is governed by a velocity potential $\phi$. The velocity components of the fluid $u$ in the $x$-direction and $v$ in the $y$-direction, are given by $\langle u,v\rangle=\nabla \phi$. Find the velocity components associated with the velocity potential $\phi(x,y)=\sin(\pi x)\sin(2\pi y)$.
\end{enumerate}

\end{document}
